\section{Notação, Estruturas de Dados e Fluxo de Execução}
\label{sec:notacao_fluxo}

\subsection{Convenções matemáticas e computacionais}

Para manter consistência entre texto, pseudocódigo e implementação em C++, este trabalho adota um padrão único de notação.
A Tabela~\ref{tab:padrao_notacao} resume os símbolos mais usados e seu mapeamento direto para o código.

\begin{table}[htbp]
\centering
\caption{Padrão de notação e mapeamento para a implementação.}
\label{tab:padrao_notacao}
\begin{tabular}{lll}
\hline
Símbolo & Significado & Nome no código/CLI \\
\hline
$S$ & solução completa & \texttt{Solution} \\
$r$ & rota individual & \texttt{Route} \\
$NR$ & clientes não roteirizados & \texttt{unrouted} \\
$d_{ij}$ & distância/custo entre $i$ e $j$ & \texttt{inst.dist[i][j]} \\
$Q$ & capacidade do veículo & \texttt{inst.capacity} \\
$m$ & limite de veículos & \texttt{inst.vehicleCount} \\
$\varepsilon$ & tolerância numérica & \texttt{eps} \\
$\mu$ & peso do termo de inserção espacial & \texttt{mu} \\
$N_{\max}$ & limite de iterações & \texttt{maxIters} \\
$L$ & limite sem melhoria & \texttt{maxNoImprove} \\
$\theta$ & tenure tabu & \texttt{tenure} \\
$\alpha$ & parâmetro RCL do GRASP & \texttt{alpha} \\
$\gamma$ & intensidade reativa do GRASP & \texttt{reactiveGamma} \\
\hline
\end{tabular}
\end{table}

\paragraph{Critério de comparação de soluções.}
No código, a comparação entre duas soluções segue regra lexicográfica custo-primeiro:
\begin{enumerate}
  \item menor custo total;
  \item em empate numérico (dentro de $\varepsilon$), menor número de rotas.
\end{enumerate}
Essa política está implementada em \texttt{Solution::betterCostThenRoutes} e é usada de forma uniforme em VND, GRASP e Tabu.

\subsection{Estruturas de dados centrais}

A implementação é organizada em três estruturas principais (\texttt{Instance}, \texttt{Route} e \texttt{Solution}), mas o ponto importante para reprodutibilidade é o \textbf{tipo concreto} das estruturas.
A Tabela~\ref{tab:estruturas_cpp} explicita se cada componente é vetor, matriz ou conjunto.

\begin{table}[htbp]
\centering
\caption{Estruturas de dados concretas da implementação (C++).}
\label{tab:estruturas_cpp}
\begin{tabular}{lll}
\hline
Componente & Tipo C++ & Papel \\
\hline
Clientes da instância & \texttt{std::vector<Customer>} & vetor indexado por id \\
Distâncias $d_{ij}$ & \texttt{std::vector<std::vector<double>>} & matriz densa $n \times n$ \\
Vizinhos próximos & \texttt{std::vector<std::vector<int>>} & lista k-NN por cliente \\
Filtro de vizinho & \texttt{std::vector<std::unordered\_set<int>>} & consulta $O(1)$ de pertinência \\
Sequência da rota & \texttt{std::vector<int>} & ordem dos clientes (inclui depósito) \\
Conjunto de rotas & \texttt{std::vector<Route>} & solução completa \\
Lista tabu & \texttt{std::vector<int>} & iteração-limite por cliente \\
Estatísticas VND & \texttt{std::array<...,4>} & contador fixo por vizinhança \\
\hline
\end{tabular}
\end{table}

\paragraph{Leitura operacional.}
No código, a estrutura mais custosa em memória é a matriz de distâncias (\texttt{dist}), pois armazena todos os pares de clientes.
Essa escolha é intencional: evita recomputar distância durante inserções e movimentos de busca local.
Já a sequência de cada rota é um vetor linear (\texttt{seq}), o que simplifica operações de inserção, remoção e troca de posição usadas por I1, VND, GRASP e Tabu.

A rotina \texttt{Route::recompute} é o núcleo de factibilidade local: percorre o vetor \texttt{seq}, atualiza tempo de chegada/início de serviço/saída, valida janela de tempo e capacidade, e recalcula custo.
A rotina \texttt{Solution::recompute} remove rotas triviais \texttt{[0,0]}, recompõe cada rota e soma o custo global.

\paragraph{Validação final.}
Depois de cada execução do solver, a função \texttt{validateSolution} faz checagem global completa:
\begin{itemize}
  \item número de rotas não excede $m$;
  \item toda rota começa e termina no depósito;
  \item nenhum cliente é repetido e nenhum cliente fica sem atendimento;
  \item todas as rotas permanecem factíveis em capacidade e janelas.
\end{itemize}

\subsection{Estruturas auxiliares por algoritmo}

Além das estruturas básicas, cada método usa estruturas específicas:
\begin{itemize}
  \item \textbf{I1}: \texttt{BestInsertion} para armazenar melhor posição por cliente e os valores $c_1,c_2$;
  \item \textbf{VND}: \texttt{VndRunStats} para contar tentativas, aplicações e ganho acumulado por vizinhança;
  \item \textbf{GRASP}: \texttt{InsertionCandidate} e vetores de probabilidades/contadores para a versão reativa;
  \item \textbf{Tabu}: \texttt{Move} (atributos do movimento) e vetor \texttt{tabuUntil[c]} para memória de curto prazo.
\end{itemize}

\subsection{Fluxo geral do solver}

O fluxo implementado no comando \texttt{vrptw solve} é mostrado no Algoritmo~\ref{alg:fluxo_global_solver}.

\begin{algorithm}[htbp]
\caption{Fluxo global do solver}
\label{alg:fluxo_global_solver}
\begin{algorithmic}[1]
\Require algoritmo \texttt{algo}, instância e parâmetros de execução
\Ensure solução final validada e artefatos de saída
\State lê instância VRPTW e BKS (quando disponível)
\State cria diretórios de logs/soluções
\State impõe $\mu=1$ para manter comparação alinhada
\If{\texttt{algo = i1}}
  \State $S \leftarrow I1$
\Else
  \If{\texttt{algo = vnd}}
    \State $S \leftarrow I1$; $S \leftarrow VND(S)$
  \Else
    \If{\texttt{algo = grasp}}
      \State $S \leftarrow GRASP$ (fixo ou reativo), com refinamento VND em cada iteração
    \Else
      \If{\texttt{algo = tabu}}
        \State $S \leftarrow$ Busca Tabu clássica inicializada em I1
      \Else
        \State gera erro de parâmetro inválido
      \EndIf
    \EndIf
  \EndIf
\EndIf
\State valida factibilidade global de $S$
\State grava solução \texttt{.sol}, log de execução e logs de convergência (quando aplicável)
\State \Return $S$
\end{algorithmic}
\end{algorithm}

\paragraph{Texto complementar ao pseudocódigo.}
Há dois pontos importantes no fluxo real.
Primeiro, todos os algoritmos passam pela mesma validação final antes da escrita da solução, então resultados inviáveis são bloqueados imediatamente.
Segundo, o formato de log é unificado, com campos de custo, gap, tempo, número de rotas e parâmetros efetivos; isso simplifica a consolidação estatística e reduz risco de inconsistência na comparação entre métodos.
